% The pism-terra model description paper, on the Copernicus template.
%
% The text of each section is written in sections/<name>.md and converted by
% the Makefile into build/<name>.tex, which this file inputs. Keep journal
% structure here (title, authors, the Copernicus section commands, the
% availability statements) and prose in the Markdown files.

\documentclass[gmd, manuscript]{copernicus}

\graphicspath{{figures/}{../docs/source/paper/figures/}}

\begin{document}

\title{pism-terra v1.0: automated simulations of any glacier and ice sheet with the Parallel Ice Sheet Model}

\Author[1]{Andy}{Aschwanden}

\affil[1]{Geophysical Institute, University of Alaska Fairbanks, Fairbanks, AK, USA}

\correspondence{Andy Aschwanden (aaschwanden@alaska.edu)}

\runningtitle{pism-terra v1.0}
\runningauthor{Aschwanden et al.}

\received{}
\pubdiscuss{}
\revised{}
\accepted{}
\published{}

\firstpage{1}

\maketitle

\begin{abstract}
\input{build/abstract}
\end{abstract}

\introduction
\input{build/introduction}

\section{Background}
\input{build/background}

\section{Case studies}
\input{build/case_studies}

\conclusions
\input{build/conclusions}

\codeavailability{TEXT}

\dataavailability{TEXT}

\authorcontribution{TEXT}

\competinginterests{The authors declare that they have no conflict of interest.}

\begin{acknowledgements}
TEXT
\end{acknowledgements}

\bibliographystyle{copernicus}
% refs: pism-terra's (docs/source/refs.bib); ice-bib: PISM's, fetched by the Makefile.
\bibliography{refs,ice-bib}

\end{document}
